\item \textbf{Vibraciones en moléculas diatómicas y potencial de Lennard-Jones.}
Asume una ayudantía de investigación con una física molecular. Ella está estudiando las vibraciones de las moléculas diatómicas. En estas vibraciones, los dos átomos de la molécula se mueven hacia adelante y hacia atrás a lo largo de la línea que los conecta (véase la figura 20.5c). Como introducción a su investigación, ella le pide que se familiarice con el potencial de Lennard-Jones (vea el ejemplo 7.9), que describe la función de energía potencial para una molécula diatómica.

Ella le pide que determine la constante de resorte efectiva, en términos de los parámetros $\sigma$ y $\epsilon$, para el enlace que mantiene unidos a los átomos en la molécula para pequeñas vibraciones alrededor de la separación de equilibrio requerida $r_{\text{eq}}$.
\textit{Sugerencia:} El ejemplo 7.9 proporciona la derivada de la función de energía potencial. Compare eso con la ecuación $F = -kx$ para encontrar la fuerza entre los átomos. Desea mostrar que $F$ tiene la forma $-kx$ y encontrar $k$. Sea la distancia de separación $r = r_{\text{eq}} + x$, donde $x$ es pequeño y aproveche las aproximaciones de la serie en la sección B.5 del Apéndice.